\section{Proof Sketches}
\label{app:proofs}

\subsection{Theorem~\ref{thm:risk}: Cost-Aware Excess Risk}

Condition on the grid-construction split $D^{(G)}$, so the feasible threshold set $\mathcal Q_G$ is fixed. For any threshold pair $q\in\mathcal Q_G$, the label-wise loss is bounded by $C_{\mathrm{TP}}=\max(\alpha_{\mathrm{cost}},\gamma_{\mathrm{abs}})$ on TP examples and $C_{\mathrm{FP}}=\max(\beta_{\mathrm{cost}},\gamma_{\mathrm{abs}})$ on FP examples. Applying Hoeffding's inequality separately to the TP and FP strata of $D^{(E)}$ and union-bounding over labels and $|\mathcal Q_G|$ threshold pairs yields a uniform label-wise deviation bound. The empirical-risk-minimization decomposition
\[
R(\hat q)-R(q_G^\star)=
\bigl(R(\hat q)-\hat R_E(\hat q)\bigr)+
\bigl(\hat R_E(\hat q)-\hat R_E(q_G^\star)\bigr)+
\bigl(\hat R_E(q_G^\star)-R(q_G^\star)\bigr)
\]
has a non-positive middle term by definition of $\hat q$. Controlling the two outer terms with the uniform deviation event gives Theorem~\ref{thm:risk}. The result is label-specific because the concentration rates depend on $m^{(E)}_{\mathrm{TP}}$ and $m^{(E)}_{\mathrm{FP}}$ separately.

\subsection{Theorem~\ref{thm:conformal}: Split-Conformal Validity}

Let $s_1,\ldots,s_m$ be calibration nonconformity scores and $s_{m+1}$ the score of a test alert under the true label. Under exchangeability, the rank of $s_{m+1}$ among the $m+1$ scores is uniform~\cite{vovk2005}. Selecting the usual finite-sample conformal quantile excludes the true label with probability at most $\alpha_{\mathrm{conf}}$, so $\Pr[Y\in\hat C(X)]\ge 1-\alpha_{\mathrm{conf}}$. EVICT then maps $|\hat C(X)|=1$ to an accepted TP/FP decision and all other set sizes to abstention.

\subsection{Theorem~\ref{thm:label_conditional}: Label-Conditional Validity}

Partition the calibration set by label. Conditional on the event that the test label is TP, the TP-labeled calibration scores and test score are exchangeable within the TP stratum; the same rank argument yields TP coverage at least $1-\alpha_{\mathrm{TP}}$. The FP case is identical. This proves separate class-conditional coverage and justifies using $\alpha_{\mathrm{TP}}<\alpha_{\mathrm{FP}}$ when false negatives are more costly than false positives.

\subsection{Theorem~\ref{thm:exact_opt}: Exact Threshold Search}

Sort the TP and FP nonconformity scores. Any label-conditional threshold can be represented by an index into the sorted TP list and an index into the sorted FP list, with an additional index for $\infty$. Therefore the finite candidate set has size at most $(m_{\mathrm{TP}}+1)(m_{\mathrm{FP}}+1)$. Each pair induces deterministic counts of correct singleton predictions, wrong singleton predictions, empty-set abstentions, and multi-label abstentions. Exhaustively evaluating empirical selective risk over the grid returns a global minimizer in $O(m^2)$ time after sorting.

\section{Algorithms}
\label{app:algorithm_details}

\subsection{Algorithm A1: EvidencePack Construction}

\begin{algorithm}[h]
\algcaption{EvidencePack Construction}
\label{alg:evidence}
\textbf{Input:} Static-analysis alert $a$ and SARIF output $O$.\\
\textbf{Output:} EvidencePack $E=(\text{slice},\text{flow},\text{constraints},\text{metadata},\text{flags})$.
\begin{enumerate}
\item Parse SARIF to recover analyzer, rule ID, CWE, severity, primary location, and source/sink locations when available.
\item Extract the analyzer-reported flow path and call chain; if no path exists, set \texttt{flow\_partial}.
\item Compute bounded backward and forward slices around the sink and source.
\item Normalize branch conditions and path predicates from the slice AST.
\item Record missing callees, unresolved reflection/native calls, missing paths, and parser failures as explicit completeness flags.
\item Return the structured EvidencePack.
\end{enumerate}
\end{algorithm}

\subsection{Algorithm A2: LLM Verification}

\begin{algorithm}[h]
\algcaption{Schema-Guided LLM Verification}
\label{alg:verify}
\textbf{Input:} Alert $a$, EvidencePack $E$, prompt mode $m\in\{\text{base},\text{contrastive}\}$, model $\mathcal L$.\\
\textbf{Output:} Label proposal $\hat y\in\{\mathrm{TP},\mathrm{FP},\mathrm{ABSTAIN}\}$ and rationale.
\begin{enumerate}
\item Render the base or contrastive prompt using $a$ and $E$.
\item Ask the model to restate the analyzer claim and required bug preconditions.
\item Ask the model to evaluate each precondition against cited evidence lines.
\item If $m=\text{contrastive}$, require separate TP and FP arguments before the final decision.
\item Emit TP, FP, or ABSTAIN with confidence and an evidence-linked rationale.
\end{enumerate}
\end{algorithm}

\subsection{Algorithm A3: Conformal Calibration}

\begin{algorithm}[h]
\algcaption{Split-Conformal Singleton-Reject Calibration}
\label{alg:conformal}
\textbf{Input:} Calibration set $\mathcal D_{\mathrm{cal}}$, nonconformity score $A(x,y)$, miscoverage $\alpha_{\mathrm{conf}}$.\\
\textbf{Output:} Prediction set function $\hat C(x)$.
\begin{enumerate}
\item For every $(x_i,y_i)\in\mathcal D_{\mathrm{cal}}$, compute $s_i=A(x_i,y_i)$.
\item Set $\hat q$ to the finite-sample $(1-\alpha_{\mathrm{conf}})$ conformal quantile of $\{s_i\}_{i=1}^m$.
\item For a new alert $x$, compute $A(x,\mathrm{TP})$ and $A(x,\mathrm{FP})$.
\item Return $\hat C(x)=\{y:A(x,y)\le\hat q\}$.
\item If $|\hat C(x)|=1$, accept the singleton label; otherwise abstain.
\end{enumerate}
\end{algorithm}

\subsection{Algorithm A4: Conditional Symbolic Escalation}

\begin{algorithm}[h]
\algcaption{Conditional Symbolic Escalation}
\label{alg:symbolic}
\textbf{Input:} Abstained alert $a$, EvidencePack $E$, solver budget.\\
\textbf{Output:} Supported TP/FP label or continued ABSTAIN.
\begin{enumerate}
\item Extract path constraints and sanitizer checks from $E$.
\item Run Z3 with a 2s timeout for feasible path constraints when encodable.
\item For Java, run Java PathFinder when project harnesses are available; for C/C++, run KLEE when build artifacts are available.
\item If a feasible counterexample reaches the sink, return TP with witness.
\item If infeasibility is proved, return FP with certificate.
\item If the solver returns \texttt{UNKNOWN}, times out, or lacks constraints, return ABSTAIN.
\end{enumerate}
\end{algorithm}

\section{Prompt Templates}
\label{app:prompts}

\subsection{Base Schema Prompt}

\begin{quote}\small
You are reviewing a static-analysis alert from \texttt{\{ANALYZER\}}.

\textbf{Alert:} Rule \texttt{\{RULE\_ID\}} (\texttt{\{CWE\}}), source \texttt{\{SOURCE\_LOC\}}, sink \texttt{\{SINK\_LOC\}}.\\
\textbf{EvidencePack:} \texttt{\{EVIDENCE\_PACK\}}.\\
\textbf{Completeness flags:} \texttt{\{FLAGS\}}.

Follow exactly four steps:
\begin{enumerate}
\item Restate the analyzer claim and what would make it a real vulnerability.
\item Enumerate every precondition required for the alert to be a true positive.
\item For each precondition, cite evidence lines and mark it SUPPORTED, CONTRADICTED, or AMBIGUOUS.
\item Output exactly one decision: TP, FP, or ABSTAIN; then give confidence in $[0,1]$ and one sentence of rationale.
\end{enumerate}
If evidence is incomplete, treat affected preconditions as AMBIGUOUS and prefer ABSTAIN unless other evidence is conclusive.
\end{quote}

\subsection{Contrastive Prompt}

\begin{quote}\small
You are reviewing a static-analysis alert. Do not decide immediately.

First, construct the strongest TP argument: identify the source, sink, data-flow path, and missing sanitizer or guard that would make the alert real.

Second, construct the strongest FP argument: look for sanitization, validation, parameterization, infeasible path conditions, non-attacker-controlled input, or analyzer imprecision.

Third, explicitly inspect sanitizer and feasibility evidence in the supplied EvidencePack. Mark each relevant fact SUPPORTED, CONTRADICTED, or AMBIGUOUS.

Fourth, weigh the TP and FP arguments. Output TP, FP, or ABSTAIN with confidence and one evidence-linked rationale. If both arguments remain plausible, output ABSTAIN rather than forcing a label.
\end{quote}

\subsection{Worked CWE-89 Example}

\textbf{Input alert:} SQL injection at \texttt{UserService.java:47--50}, source \texttt{request.getParameter("id")} and sink \texttt{stmt.execute(query)}.

\begin{verbatim}
45: String id = request.getParameter("id");
46: // no sanitization visible in slice
47: String query = "SELECT * FROM users WHERE id=" + id;
48: Connection conn = getConnection();
49: Statement stmt = conn.createStatement();
50: stmt.execute(query);
\end{verbatim}

\textbf{Expected output:}
\begin{verbatim}
1. CLAIM: The analyzer reports untrusted HTTP input concatenated into
   a SQL string and executed.
2. PRECONDITIONS: P1 user-controlled input; P2 flow into SQL string;
   P3 no sanitizer or parameterized query; P4 execution by SQL sink.
3. EVIDENCE: P1 SUPPORTED line 45; P2 SUPPORTED line 47;
   P3 SUPPORTED lines 45-50 show no sanitizer or PreparedStatement;
   P4 SUPPORTED line 50.
4. DECISION: TP. CONFIDENCE: 0.94. RATIONALE: All required SQL
   injection preconditions are directly supported by the slice.
\end{verbatim}

\subsection{Self-Consistency Protocol}

For API-only models, the verifier is queried $k_{\mathrm{sc}}$ times with temperature 0.7. The main conformal studies use $k_{\mathrm{sc}}=5$; cross-validation uses $F=5$ folds. Vote share $v(y\mid x)$ is the fraction of samples selecting label $y$. Sensitivity from the provided study export: abstention is 14.9\% at $k_{\mathrm{sc}}=3$, 12.7\% at $k_{\mathrm{sc}}=5$, 11.8\% at $k_{\mathrm{sc}}=10$, and 11.3\% at $k_{\mathrm{sc}}=20$; ECE varies by $\pm 0.01$ across this range.

\section{Extended Experimental Diagnostics}
\label{app:experiment_details}

\subsection{Prediction-Set Frequencies}

The earlier 1,000-alert SpotBugs Juliet study reports 87.3\% singleton sets, 11.6\% multi-label sets, and 1.1\% empty sets overall. Per-CWE singleton frequencies are 91.4\% for CWE-89, 89.1\% for CWE-78, and 81.2\% for CWE-190. Per-CWE multi-label and empty-set frequencies were not provided in the prompt and should be filled from the artifact export before camera-ready: [VERIFY: per-CWE multi-label and empty prediction-set frequencies].

\subsection{Per-CWE Difficulty}

The earlier SpotBugs study reports pre-escalation precision of 93.1\% for CWE-89, 91.8\% for CWE-78, and 83.4\% for CWE-190. These numbers are source-specific diagnostics and are not part of the Study~2 six-model conformal table. Integer overflow is harder in that study, consistent with the need for arithmetic-range reasoning beyond simple taint flow. Per-CWE abstention and ECE values were not provided in the prompt: [VERIFY: per-CWE abstention and ECE from earlier SpotBugs export].

\subsection{Symbolic Verification Impact}

\begin{table}[h]
\centering
\caption{Symbolic characterization from the 1,000-alert study. Pre-escalation abstentions and high-severity dismissal audits are routed to post-escalation symbolic checks.}
\label{tab:symbolic_appendix}
\small
\begin{tabular}{lc}
\toprule
\textbf{Metric} & \textbf{Value} \\
\midrule
Symbolic invocations & 184/1000 (18.4\%) \\
From pre-escalation abstention & 127/184 \\
From high-severity dismissal audit & 57/184 \\
SMT \texttt{UNKNOWN} / timeout retained as abstain & 15/184 (8.1\%) \\
LLM errors corrected post-escalation & 23 \\
FP$\rightarrow$TP safety-critical corrections & 17 \\
TP$\rightarrow$FP efficiency corrections & 6 \\
Average symbolic verification time & 4.2s \\
Z3 timeout & 2s \\
JPF/KLEE timeout & 10s \\
\bottomrule
\end{tabular}
\end{table}

\subsection{Cost and Latency}

\begin{table}[h]
\centering
\caption{Operational cost from the provided cost study. Risk--coverage coordinates were not included in the prompt and remain to be verified from the artifact export.}
\label{tab:cost_appendix}
\small
\begin{tabular}{lccc}
\toprule
\textbf{Method} & \textbf{Tokens/alert} & \textbf{Latency/alert} & \textbf{Solver timeout rate} \\
\midrule
Evidence-free prompting & 820 & 1.1s & 0 \\
EVICT full pipeline & 1,710 & 2.7s & 8.1\% \\
Risk--coverage curve & [VERIFY: tokens by coverage] & [VERIFY: latency by coverage] & [VERIFY: solver utilization] \\
\bottomrule
\end{tabular}
\end{table}

\subsection{False-Negative and Cost Sensitivity}

The provided prompt requests false-negative analysis under abstention and sensitivity to $\alpha_{\mathrm{cost}}/\beta_{\mathrm{cost}}/\gamma_{\mathrm{abs}}$, but does not provide numeric outputs. The review draft therefore marks these as artifact-fill items rather than fabricating values: [VERIFY: abstention rate on true vulnerabilities by study], [VERIFY: false-negative rate among auto-dismissed alerts], and [VERIFY: cost-sensitivity sweep over $\alpha_{\mathrm{cost}}$, $\beta_{\mathrm{cost}}$, and $\gamma_{\mathrm{abs}}$].

\section{Glossary}
\label{app:glossary}

\begin{description}
\item[ECE] Expected calibration error; a binned difference between model confidence and empirical accuracy.
\item[TP/FP/FN/TN] True positive, false positive, false negative, and true negative under alert validity labels.
\item[$\hat q$] The conformal nonconformity threshold estimated from calibration data.
\item[SARIF] Static Analysis Results Interchange Format, a standard JSON format for analyzer outputs.
\item[EvidencePack] EVICT's structured alert record containing code slice, data-flow summary, path constraints, metadata, and completeness flags.
\item[Singleton prediction set] A conformal set containing exactly one label; EVICT accepts it as TP or FP.
\item[Multi-label prediction set] A conformal set containing both TP and FP; EVICT abstains because both labels remain plausible.
\item[Empty prediction set] A conformal set containing no labels; EVICT abstains because neither label is sufficiently supported.
\item[Exchangeability] The calibration and test examples are identically distributed up to permutation, the standard assumption behind split-conformal validity.
\item[Vote share] The fraction of self-consistency samples choosing a label.
\item[ECE bins] Confidence intervals used to aggregate calibration error; binning choices affect the diagnostic but not conformal validity.
\end{description}

\section{Reproducibility Statement}
\label{app:reproducibility}

For double-blind review, the artifact is described anonymously. We will release the pipeline implementation, preprocessing scripts, prompts, model-version metadata, inference settings, calibration folds, and evaluation logs. The release will include the unrun RQ5 frontier-model harness and projection method so that groups with sufficient API budget can execute the frontier evaluation. The review-mode paper avoids repository names, local paths, author names, affiliations, acknowledgments, and grant identifiers; these can be restored only in a de-anonymized camera-ready version.
